\chapter{Происхождение горизонтальной плоской моды}
\label{ch:v8-horizontal-flat-direction-parent-lift-gate}

Единственная некалибровочная нулевая мода quotient-гессиана допускает
точную изотипическую интерпретацию. Вакуумная инцидентность распадается на
три стрелки целевого $up$-блока и семь остальных стрелок. Их независимые
фазовые касательные $z_u,z_r$ принадлежат полному transfer-носителю и имеют
следовую метрику
\begin{equation}
 Z=(z_u,z_r),
 \qquad Z^{\mathsf T}K_{\Phi}Z=\operatorname{diag}(6,20),
 \qquad H_{\Phi}Z=0.
 \label{eq:v8-horizontal-phase-plane}
\end{equation}
Связь этой фазовой плоскости с независимым репером нарушенной
калибровочной орбиты равна
\begin{equation}
 X^{\mathsf T}K_{\Phi}Z=
 \begin{pmatrix}0&0\\0&0\\-6&8\end{pmatrix}.
 \label{eq:v8-horizontal-phase-orbit-coupling}
\end{equation}
Поэтому единственная примитивная горизонтальная комбинация имеет
коэффициенты
\begin{equation}
 v_0=4z_u+3z_r,
 \qquad X^{\mathsf T}K_{\Phi}v_0=0,
 \qquad v_0^{\mathsf T}K_{\Phi}v_0=276.
 \label{eq:v8-horizontal-phase-four-three}
\end{equation}

Это не случайность квадратичного приближения. Для $z\in U(1)$ положим
\begin{equation}
 U(z)=\operatorname{diag}(z^4I_3,z^3I_7),
 \qquad A(z)=U(z)A_0.
 \label{eq:v8-horizontal-phase-circle}
\end{equation}
Из структуры инцидентности следует точное тождество
\begin{equation}
 A(z)A(z)^*=A_0A_0^*,
 \qquad A(z)^*A(z)=A_0^*A_0.
 \label{eq:v8-horizontal-phase-gram-invariance}
\end{equation}
Следовательно, любой родитель, построенный только из грамовых концов и
следов их слов, постоянен на всей этой окружности и не способен поднять
$v_0$ ни на каком порядке.

Обычный determinant также не возникает автоматически: для прямоугольника
$A_0:\mathbb C^{11}\to\mathbb C^{10}$ максимальные миноры образуют
\begin{equation}
 \dim\!\left(\Lambda^{10}\mathbb C^{11}\right)^*
 \dim\!\left(\Lambda^{10}\mathbb C^{10}\right)
 =\binom{11}{10}\binom{10}{10}=11,
 \label{eq:v8-horizontal-maximal-minor-carrier}
\end{equation}
то есть ковектор determinant-line, а не канонический скаляр. Для
фазочувствительного члена требуется дополнительное правило свёртки или
ориентация, отсутствующие в текущем родителе.

\begin{theorem}[Точная природа плоской моды]
Горизонтальная нулевая мода является остаточной изотипической фазой с
отношением весов $4:3$. Она сохраняет оба грамовых конца точно на всей
$U(1)$-окружности.
\end{theorem}

\begin{nogocriterion}[Запрет грамового родительского подъёма]
Никакая функция от $A A^*$, $A^*A$ и их следовых слов не создаёт массу
для $v_0$. Прямоугольный maximal-minor объект имеет размерность $11$ и не
становится скаляром без нового contraction/orientation datum. Добавление
ручного массового члена запрещено.
\end{nogocriterion}

\begin{protocol}
\begin{enumerate}
 \item фазовая плоскость: \textbf{$2$};
 \item её метрика: \textbf{$\operatorname{diag}(6,20)$};
 \item горизонтальная комбинация: \textbf{$4:3$};
 \item точная грамова симметрия: \textbf{сохраняется на всех порядках};
 \item грамовый родитель поднимает моду: \textbf{нет};
 \item канонический скалярный determinant: \textbf{нет};
 \item следующий гейт: \textbf{допуск фазочувствительной determinant-line
 свёртки}.
\end{enumerate}
\end{protocol}

Машинный сертификат сохранён в
\path{s2t_v8_horizontal_flat_direction_parent_lift_gate_results.json}.